\chapter{Discussion}
\label{ch:discussion}
% Red \resultpending{...} markers are numbers from the held-out classroom test run
% (chapters 7--10); fill them from those chapters once the run is recorded in the ledger.
\providecommand{\resultpending}[1]{\textcolor{red}{[\,#1\,]}}

This chapter reads the results of \cref{ch:rq1,ch:rq2,ch:rq3,ch:rq4} together,
asks what they mean for programming assessment, and sets out what the evidence
cannot show. Development figures come from the dev split of the pilot
(\cref{sec:split}) and are marked as such; headline figures come from the single
held-out test run.

\section{Summary of findings}
\label{sec:disc-findings}

\paragraph{RQ1: a mid-session switch can be detected, slowly, at a real cost in
false alarms.} On held-out classroom data the frozen detector catches
\resultpending{RQ1 test: \% of same-device splices caught [95\% CI]} of spliced
takeovers with a median delay of \resultpending{RQ1 test: median delay}, at
\resultpending{RQ1 test: false alarms per active hour [CI]} on real sessions. The
development split gave 64\% (49--78\%) at a median of 450\,s and 1.59 false
alarms per active hour (\cref{tab:retune-splice}). On simulated typists the same
design caught 86\% of speed-matched takeovers at 0.19 false alarms per hour
(\cref{sec:rq1-harness}). The gap between those two settings is itself a
finding (\cref{sec:disc-drift}). Among the compared methods,
\resultpending{RQ1 test: best and worst baseline vs.\ thesis detector on splices}.

\paragraph{RQ2: transcription measures and copying signs.}
\resultpending{RQ2 test: agreement of L-scale and process score with copying
signs (AUC or rank correlation, CI)}. On the pilot, most pastes were the
student's own code moving inside the editor (434 of 518 on dev, 84\%), and
sessions with an outside paste did not raise more detector episodes than other
sessions (\cref{sec:retune}). Copying is therefore reported as its own signal
rather than folded into the switch detector.

\paragraph{RQ3: typing moves between devices.} The first look found 9.9
episodes per active hour on a device other than the student's usual one against
5.7 on the usual one (\cref{tab:pilot}), and a robust rhythm distance of 2.1
between two sessions of the same student against 2.9 between two students.
\resultpending{RQ3 test: within-student drift same vs.\ other device class, and
false alarms with a cross-device vs.\ same-device baseline}.

\paragraph{RQ4: precision is low at realistic rates.} At 1 in 100 sessions,
\resultpending{RQ4 test: precision at 1/100 [CI]} of flags would be real; at 1 in
500, \resultpending{RQ4 test: precision at 1/500 [CI]}. As an illustration from
the dev split: with 64\% of splices caught and 29\% of real sessions flagged at
least once, precision would be about 2\% at 1 in 100 and 0.4\% at 1 in 500, so a
teacher would see around 45 honest sessions flagged for each real takeover at
the higher rate. Subgroup differences:
\resultpending{RQ4 test: largest subgroup gap in false alarms (group, rate, CI)}.

\section{What the results mean for assessment}
\label{sec:disc-assessment}

\subsection{Real classwork drifts more than any simulation}
\label{sec:disc-drift}
The most consistent result of the thesis is that real classroom sessions look
less like their own baseline than simulated sessions do. On synthetic typists
the detector met its target of 0.2 false alarms per hour; the first look at real
labs gave 7.5 episodes per active hour, and even after the pilot fixes the dev
split gives 1.59. Re-tuning located the cause (\cref{sec:retune}): the
standardised score of honest work averages 1.2--1.6 rather than 0 on every
channel. A student's baseline is a few short sessions on one or two questions,
early in the term; later sessions are on harder questions, at another time of
day, sometimes on another machine, by a student who is still learning to
program and to type. None of that is a change of author, and all of it moves the
features.

This agrees with the critiques of behavioural-biometric evaluation reviewed in
\cref{sec:rw-keystroke}: drift between sessions, not separability within one,
limits continuous authentication \cite{eberz2017evaluating}, and natural typing
on varied hardware is harder than prompted typing
\cite{sun2016buffalo,murphy2017clarkson}. It also echoes writing research, where
the task changes keystroke features \cite{conijn2019keystroke}. For this
thesis, it means that synthetic results are a test of the code, not of the
method, and that any threshold must be calibrated on classroom data from the
course where it will be used.

\subsection{No operating point is free}
Every change that cut false alarms on real sessions also cut detection of
spliced takeovers (\cref{tab:retune-splice}): dropping build order, raising
$k$ to 1, 1.5 and 2 traced a curve from 64\% caught at 1.59 false alarms per hour
to 11\% at 0.04. The quiet configuration that looked best from flag rates alone
hid the takeovers along with the drift. The thesis therefore reports the whole
trade-off curve rather than a single tuned point, and the choice of operating
point is left to the people who will act on the flags. For a tool whose flags
lead only to a conversation, the frozen configuration (more flags, more
takeovers caught) was the one chosen for the live system.

\subsection{Which signals carry the evidence}
Three channels contribute differently. The rhythm-based identity channel behaves
on real data much as it did on simulated data (0.24 episodes per active hour in
the first look). The key-pair channel fires first on most of the caught dev
splices, which makes its vocabulary on real classes the most promising place to
gain detection without more false alarms. The process channel (pauses,
revisions, build order, runs, focus) raised most of the pilot's flags, because
build order follows the question as much as the person; yet it also carries part
of the splice detection, partly because a spliced partner continues their own
code (\cref{sec:disc-limitations}). The identity channels are therefore also
reported on their own in \cref{ch:rq1}:
\resultpending{RQ1 test: identity channels alone, \% caught and false alarms per hour}.

\subsection{Takeovers are caught in minutes, not seconds}
A median delay of several minutes (450\,s on dev) means a helper who types for
two or three minutes will usually go unflagged, while one who takes over for the
rest of a lab is likely to be caught, with an estimated switch time the teacher
can check on the timeline. The tool therefore deters and surfaces sustained
takeovers, the case that matters most for assessment, rather than every
borrowed keystroke. A full proxy who types the whole session has no change to
detect within that session; only a comparison with the student's earlier
sessions (the baseline) can reveal it, and only if the baseline itself was the
student's own work.

\subsection{Copying needs its own evidence}
Most pastes in real classes were the student's own text, and outside pastes did
not predict detector episodes. Paste origin and blocked outside pastes
(\cref{sec:paste-guard}) are direct, explainable evidence of copying, far more
so than a shift in typing rhythm. This matches the wider finding that detecting
generated text after the fact is unreliable and can be unfair
\cite{weberwulff2023testing,liang2023gptdetectors}: what the editor can record
about \emph{how} text arrived is more defensible than any guess about where it
came from.

\section{A flag is a conversation}
\label{sec:disc-conversation}
The precision figures settle what a flag may be used for. At any plausible
prevalence most flags are honest students, so a flag cannot be evidence of
misconduct; at most it is a reason to look and to talk. The design commits to
this throughout: students are told on the consent screen that a flag only ever
leads to a conversation, the review timeline shows the evidence behind each
moment (which channel, which features, the estimated switch time, the code
written around it) rather than a verdict, and nothing feeds a grade
(\cref{sec:review-timeline}).

How a teacher uses a moment matters as much as how it was computed. A teacher
who opens a moment should first look at the code and the timeline, and then talk
with the student the way a lab viva already works: ask them to explain a part of
the code written after the estimated switch, to extend it, or to fix a small bug
in it. A student who wrote it can usually do so, whatever their typing looked
like. The conversation should not open with the flag as an accusation, and its
outcome should be recorded as what it was (\emph{explained, own work},
\emph{had help}, \emph{copied}, \emph{not discussed}) and nothing more
(\cref{sec:teacher-review}). On the dev split, teachers recorded
\resultpending{conversation outcomes: counts by outcome (tab:teacher-outcomes)};
these describe how moments were used, not whether they were right.

False flags carry harms even under this design. A student asked about their work
may feel distrusted, as students monitored by proctoring software report
\cite{balash2021examining,coghlan2021proctor}; repeated flags for the same
student could shape how a teacher sees them; and flags do not fall evenly. The
pilot already shows uneven rates by device (9.9 against 5.7 episodes per hour
on a less usual device), and a student on a phone has no typing rhythm at all.
Students who own no laptop, or who switch between a shared lab PC and a borrowed
one, would be flagged more for reasons unrelated to honesty. The mitigations are
procedural: labs should require a laptop or PC, teachers should know that device
changes raise flags, and a flag should never be raised with a student twice for
the same pattern without a new reason. Prinsloo and Slade's obligation to act
\cite{prinsloo2017elephant} cuts both ways here: having the data creates a duty
to use it well, which includes the duty not to over-read it.

Keeping a human decision at every step is also what the EU AI Act requires of
high-risk systems in education (human oversight, transparency to the people
affected) \cite{euaiact2024}, and what India's DPDP Act's purpose limitation
implies: data collected to support teaching is used for a conversation about
learning, not for discipline \cite{dpdp2023}.

\section{Limitations}
\label{sec:disc-limitations}

\paragraph{No ground truth.} The thesis draws every conclusion from naturalistic
classroom recordings (\cref{sec:naturalistic}). No real session is known to
contain a switch and none is known to be clean. False alarms measured on real
sessions are therefore upper bounds, since some of them may be real switches, and
precision computed from them is a lower bound. Teachers' conversation outcomes
are not labels either: they exist only for moments a teacher chose to discuss,
and a student's account is not proof.

\paragraph{Spliced takeovers approximate real ones.} A splice joins two real
students' typing on the same question at a known time (\cref{sec:splicing}). It
differs from a real takeover in ways that cut both ways. Student $B$ continues
their own code, not $A$'s, so line-level build order reflects $B$'s document,
and a real helper would first pause to read $A$'s code; both make the process
channel's job easier than in reality, so splice detection is optimistic for that
channel. On the other hand, a real helper is often a stronger programmer than
the student and types differently, which splices between classmates on the same
question do not guarantee. Splices also assume each source session had a single
author. Finally, splices need two students on the same question and split side:
45 on dev, and only 2--3 per run that pair different device classes, too few to
measure a takeover by someone on another kind of machine.

\paragraph{The device confound.} Students chose their devices. Lab PCs and
personal laptops differ in keyboard, operating system and browser, and the
students who brought laptops may differ from those who did not. Device class is
itself inferred, from a browser fingerprint shared by three or more students
(\cref{sec:baseline}), which can misplace a laptop model popular in the class.
Phones (120 sessions in the first look) send no key events, so those sessions
have no identity evidence at all. RQ3 compares devices within students only,
and RQ4 reports device class with this caveat, but a device change and an
author change remain partly confounded: both move the same features.

\paragraph{Modest sample, short sessions, short baselines.} The pilot covers
1,118 sessions of 281 students over four weeks (\cref{tab:pilot}), but the median
session is five minutes and 108 key presses, and many sessions are opened and
left (430 of 521 idle dev sessions have under 50 events). Only 332 dev sessions
had a ready baseline and some typing, 10.7 active hours, and most students'
baselines are a few minutes from one or two sessions. Confidence intervals are
correspondingly wide (64\%, 49--78\% on dev), and subgroup comparisons in
\cref{ch:rq4} can detect only large differences. Baselines come from a student's
earliest minutes, so a student who improves over the term drifts away from
their baseline for honest reasons.

\paragraph{One institution, one term, two languages.} All data comes from IIIT
Surat classes in C and Python in one term.
Typing behaviour, devices, class sizes and integrity norms differ elsewhere, and
the results may not transfer to other courses, languages or keyboard layouts.

\paragraph{Ethics approval.} The pilot ran with oral institutional approval from
the Dean, the Head of Department and the course faculty; written confirmation is
pending (\cref{sec:ethics}, \cref{app:ethics}). Students saw a consent screen
before any recording, raw key events are deleted after 28 days, the analysis
used pseudonymous aggregates only, and spliced sessions never left the analysis.
Even so, an examiner may expect written approval from an institutional ethics
committee for human-subject data, and the thesis can claim only what the oral
approval covered until that confirmation is attached.

\paragraph{Out of scope.} Deliberate imitation of another person's typing, a
helper working through remote-desktop software on the student's machine,
dictation, and help that the student then retypes (which is transcription, the
subject of \RQ{2}, not a change of typist) are not modelled. The detector also
cannot tell a takeover by a classmate from legitimate help by a teaching
assistant who typed a line; the conversation can. Demographic attributes such
as gender or prior experience were not recorded, so fairness is reported only
for groups visible in the recordings.

\section{Threats to validity}
\label{sec:disc-validity}
\begin{description}[leftmargin=1.5em,style=nextline]
  \item[Internal validity] The split by student and by week (\cref{sec:split})
    kept tuning off the test data, but the first look at the pilot reported
    aggregates over all students before the split was fixed; it informed the
    decision to re-tune, though no parameter was fitted on it. Test-split
    sessions also feed class-level parts of the live review, which the frozen
    evaluation configuration fits on dev students only.
  \item[Construct validity] A splice stands in for a takeover and copying signs
    stand in for copying; neither is the real event. Episodes per active hour
    depend on the definition of an active window (30\,s with 20 code key
    presses), which excludes idle time but also short bursts of typing.
  \item[External validity] One institution, one term, two languages and the
    device mix of one college, as above. The classes were real, which is the
    design's strength: nobody was asked to cheat, and behaviour was not changed
    by staging.
  \item[Conclusion validity] Intervals come from resampling students, so they
    account for repeated sessions of the same student, but subgroups and
    cross-device splices are small. The held-out test was scored once per
    frozen configuration and recorded in the ledger, so the reported test
    figures are not the best of several tries.
\end{description}

\section{Recommendations for practice}
\label{sec:disc-practice}
For a department that wants to use process evidence in programming labs, the
results suggest five rules. First, calibrate on your own classes: thresholds set
on simulated or other courses' data will flag several times too often. Second,
require a laptop or PC for assessed labs, and treat flags raised on a student's
less usual device with extra caution. Third, treat paste origin and blocked
outside pastes as the evidence for copying, and typing changes as the evidence
for a change of author; do not mix them. Fourth, choose the operating point
knowingly, from the whole trade-off curve, and say which one is in use. Fifth,
use every flag only to start a conversation about the code, record the outcome
as what it was, and never let a flag alone affect a grade.
